\documentclass[11pt,a4paper]{article}
\usepackage[T1]{fontenc}
\usepackage[utf8]{inputenc}
\usepackage{lmodern,microtype}
\usepackage{amsmath,amssymb,amsthm,mathtools,mathrsfs}
\usepackage[margin=27mm,headheight=14pt]{geometry}
\usepackage{booktabs,longtable,array,enumitem}
\usepackage[dvipsnames]{xcolor}
\usepackage{hyperref}
\usepackage{fancyhdr}
\hypersetup{colorlinks=true,linkcolor=MidnightBlue,citecolor=MidnightBlue,urlcolor=MidnightBlue,
 pdftitle={Surcomplex Analysis: Hahn Summation, Analytic Germs, and a Residue Calculus},
 pdfauthor={},pdfsubject={Analysis over the complexification of the surreal numbers}}
\pagestyle{fancy}\fancyhf{}
\fancyhead[L]{\small Surcomplex Analysis}
\fancyhead[R]{\small\nouppercase{\leftmark}}
\fancyfoot[C]{\thepage}
\renewcommand{\sectionmark}[1]{\markboth{\thesection.\ #1}{}}
\setlength{\parskip}{4pt}
\setlength{\parindent}{1.25em}
\setlist{topsep=4pt,itemsep=2pt,parsep=0pt}
\setlength{\emergencystretch}{2em}
\allowdisplaybreaks[2]
\newtheorem{theorem}{Theorem}[section]
\newtheorem{proposition}[theorem]{Proposition}
\newtheorem{lemma}[theorem]{Lemma}
\newtheorem{corollary}[theorem]{Corollary}
\theoremstyle{definition}
\newtheorem{definition}[theorem]{Definition}
\newtheorem{example}[theorem]{Example}
\newtheorem{construction}[theorem]{Construction}
\theoremstyle{remark}
\newtheorem{remark}[theorem]{Remark}
\newcommand{\No}{\mathbf{No}}
\newcommand{\SC}{\mathbb S}
\newcommand{\C}{\mathbb C}
\newcommand{\R}{\mathbb R}
\newcommand{\Q}{\mathbb Q}
\newcommand{\N}{\mathbb N}
\newcommand{\Z}{\mathbb Z}
\newcommand{\OO}{\mathcal O}
\newcommand{\HH}{\mathcal H}
\newcommand{\mm}{\mathfrak m}
\newcommand{\Fin}{\mathcal V}
\newcommand{\supp}{\operatorname{supp}}
\newcommand{\st}{\operatorname{st}}
\newcommand{\Res}{\operatorname{Res}}
\newcommand{\ord}{\operatorname{ord}}
\newcommand{\LC}{\operatorname{lc}}
\newcommand{\id}{\operatorname{id}}
\newcommand{\red}{\operatorname{red}}
\newcommand{\cw}{\preccurlyeq_{\mathrm{cw}}}
\newcommand{\Hs}{\mathop{\sum}\nolimits^{\!H}}
\newcommand{\HI}{\mathop{\int}\nolimits^{\!H}}
\newcommand{\Hoint}{\mathop{\oint}\nolimits^{\!H}}
\newcommand{\abs}[1]{\left|#1\right|}
\newcommand{\norm}[1]{\left\|#1\right\|}
\newcommand{\lift}[1]{#1^{\sharp}}
\newcommand{\disk}{\mathbb D}
\newcommand{\Tay}{\operatorname{Tay}}
\newcommand{\remm}{\operatorname{rem}_{<m}}
\newcommand{\dd}{\,d}
\newcommand{\ii}{i}
\title{\textbf{Surcomplex Analysis}\\[5pt]
\large Hahn Summation, Analytic Germs, and a Residue Calculus}
\author{}
\date{21 September 2026}
\begin{document}
\maketitle
\begin{abstract}
The surcomplex numbers form the algebraically closed class field
$\SC=\No[i]$. Its algebra resembles that of $\C$, but its full
order-induced topology does not support the usual sequential foundations
of complex analysis. We develop a two-level replacement. At the local
level, formal power series with surcomplex coefficients define genuine
analytic germs by Hahn summation on sufficiently small neighborhoods.
At the global level, well-ordered Hahn families of classical holomorphic
functions impose coherence between otherwise independent infinitesimal
neighborhoods. We prove Taylor and composition theorems, identity and
open-mapping theorems, a maximum-modulus principle, support-controlled
Weierstrass preparation and division, and a precise specialization theorem
for zero divisors. A coefficientwise contour integral then satisfies a
residue theorem counting actual surcomplex poles, not merely poles of
leading coefficients. Cauchy's formula, the argument principle, and
Rouch\'e-type comparison follow. We also establish Schwarz--Pick for
canonical lifts, identify the correct coefficientwise form of Liouville's
theorem, and construct explicitly different global analytic exponentials
agreeing on both the real-surreal axis and the ordinary complex plane.
All global statements specify their analytic class and their domain;
Hahn summation is never identified with convergence of partial sums.
\end{abstract}
\noindent\textbf{Mathematical status.} This is a self-contained development
of the stated framework, using standard foundations of surreal and Hahn
fields and classical one-variable complex analysis. It does not claim
priority for the constructions or a unique intrinsic analytic geometry
on all of $\No[i]$. The included computer checks concern finite symbolic
examples; they are not formal verification of the proofs.
\clearpage
\begingroup
\small
\setlength{\parskip}{0pt}
\makeatletter
\renewcommand*\l@section[2]{\addvspace{3pt}%
  \@dottedtocline{1}{0pt}{1.8em}{\bfseries #1}{\bfseries #2}}
\makeatother
\tableofcontents
\endgroup
\clearpage

\section{The problem and the choice of analytic structure}\label{sec:intro}
An algebraically closed field does not by itself determine an analysis.
The expression $\No[i]$ specifies addition, multiplication, conjugation,
and an algebraic closure; it does not specify which infinite sums,
contours, or analytic continuations are legitimate. In particular, three
operations must be separated:
\[
 \text{a classical complex sum},\qquad
 \text{a Hahn sum},\qquad
 \text{a limit in the full surreal topology}.
\]
They need not agree, and one may exist when the others do not.

The present construction retains two complementary features of ordinary
complex analysis. Formal Taylor algebra describes what happens on an
arbitrarily small scale. Classical holomorphic dependence of the Hahn
coefficients describes how different infinitesimal neighborhoods fit
together. The second requirement is essential: prescribing an unrelated
power series on each infinitesimal neighborhood produces many locally
analytic functions with no useful continuation principle.

\subsection{The two levels}
Locally, we realize every series in $\SC[[X]]$ as a germ by evaluating it
at increments smaller than all the scales appearing in its coefficients.
This gives differentiation, local inversion, ramification, and formal
residues. It also reveals why local analyticity alone is too permissive.

For a connected ordinary complex domain $U$ and a set-sized divisible
ordered subgroup $\Gamma\subseteq\No$, we then use
\begin{equation}\label{eq:H-intro}
 \HH_\Gamma(U)=
 \left\{\sum_{\gamma\in S} f_\gamma(z)t^\gamma:
 S\subseteq\Gamma\text{ well ordered},\ f_\gamma\in\OO(U)\right\}.
\end{equation}
Here $t^\gamma$ denotes the Conway monomial $\omega^{-\gamma}$.
Such a family evaluates on all finite surcomplex points whose standard
part belongs to $U$, and not just on the ordinary points of $U$.
The coefficients may have arbitrary classical growth, and $\Gamma$ may
have arbitrary ordered rank.

The principal result is the following specialization law. If
$F\ne0$ has first nonzero coefficient $f_{\gamma_0}$, then a classical
zero $a$ of $f_{\gamma_0}$ of multiplicity $m$ has exactly $m$
surcomplex zeros of $F$ in its infinitesimal neighborhood, counted with
multiplicity. This is a statement about the evaluated function. Its
proof is an explicit Hahn-supported Weierstrass preparation, not an
appeal to numerical closeness.

\subsection{Relation to existing foundations}
Conway's normal forms and real closedness provide the algebraic starting
point. Van den Dries and Ehrlich describe the surreal field as a directed
union of embedded Hahn fields and construct its restricted analytic
structure \cite{vdDE}. Berarducci and Mantova develop surreal analytic
germs and make explicit the failure of unrestricted analytic
continuation \cite{BMgerms}. Our local construction complexifies that
picture; our global choice in \eqref{eq:H-intro} enforces a particular
kind of coefficientwise holomorphic coherence.

The support calculus uses the Hahn--Neumann summability lemma
\cite{Neumann,PoonenUnits}. Algebraic closedness of the set-sized
coefficient fields follows from the standard theory of maximally
complete Hahn fields \cite{Poonen}. Classical results about holomorphic
functions are used only for the complex coefficient functions; a
standard source is \cite{McMullen}. The word ``Hahn'' is used here for
formal generalized series with constant surreal monomials. This is not
the same construction as Hahn-meromorphic functions on a logarithmic
cover of the complex plane in \cite{MS}.

\subsection{A guide to the conclusions}
\begin{center}
\small
\begin{tabular}{@{}p{0.29\textwidth}p{0.62\textwidth}@{}}
\toprule
Classical principle & Surcomplex version established here\\
\midrule
Taylor expansion and inversion & Hahn-summed analytic germs; arbitrary surcomplex coefficients.\\
Identity and open mapping & Valid for coherent Hahn families on a connected complex base.\\
Maximum modulus & A fine-topological local maximum forces a coherent family to be constant.\\
Weierstrass preparation & A holomorphic unit times a finite polynomial with infinitesimal coefficients.\\
Residue and Cauchy formulas & Coefficientwise contours recover residues at actual surcomplex poles.\\
Rouch\'e and argument principle & Multiplicity is conserved within standard-part fibers.\\
Liouville & Coefficientwise boundedness works; a single surreal bound does not.\\
Schwarz--Pick & Exact for canonical lifts of classical disk maps, not for arbitrary Hahn deformations.\\
\bottomrule
\end{tabular}
\end{center}

\section{The field, its modulus, and its small subfields}\label{sec:field}
\subsection{Normal forms and the complexification}
We work in a class theory such as NBG, or read every argument inside a
sufficiently large set-sized fragment. The class of surreal numbers is
an ordered real closed field. Consequently
\[
 \SC=\No[i]=\{x+iy:x,y\in\No\},\qquad i^2=-1,
\]
is algebraically closed. A ``polynomial over $\SC$'' always has finite
degree and a set of coefficients.

Every surcomplex number has a unique complex-coefficient normal form
\begin{equation}\label{eq:normal}
 z=\sum_{\gamma\in S}c_\gamma t^\gamma,
 \qquad c_\gamma\in\C,\quad
 t^\gamma:=\omega^{-\gamma},
\end{equation}
where $S\subseteq\No$ is a well-ordered set. Zero coefficients may be
omitted. This follows by combining the two real normal forms; a finite
union of well-ordered subsets of a linear order is well ordered.

\begin{remark}[Meaning of the exponent]
The notation $t^\gamma$ in this article is monomial notation. It obeys
$t^\gamma t^\delta=t^{\gamma+\delta}$, but is not being defined as
$\exp(-\gamma\log\omega)$ using the surreal exponential. These two
conventions for exponentiation must not be silently identified for
arbitrary surreal exponents.
\end{remark}

For $z\ne0$, set
\[
 v(z)=\min\supp z,\qquad \LC(z)=c_{v(z)},\qquad v(0)=+\infty.
\]
Then
\begin{equation}\label{eq:valuation}
 v(zw)=v(z)+v(w),\qquad
 v(z+w)\ge\min\{v(z),v(w)\}.
\end{equation}
Conjugation acts on the coefficients. The ordered-field modulus is
\[
 |x+iy|=\sqrt{x^2+y^2}\in\No_{\ge0}.
\]
It is multiplicative and satisfies the triangle inequality. The usual
algebraic proof uses only inequalities valid in a real closed field.
In particular,
\begin{equation}\label{eq:modulusleading}
 z=t^\gamma(c+\epsilon),\ v(\epsilon)>0,\ c\in\C^*
 \quad\Longrightarrow\quad
 |z|=t^\gamma(|c|+\epsilon_1),\ v(\epsilon_1)>0.
\end{equation}
Thus $v(|z|)=v(z)$.

\subsection{Finite points and infinitesimal neighborhoods}
Define
\[
 \Fin=\{z:|z|<n\text{ for some }n\in\N_{>0}\},\qquad
 \mm=\{z:|z|<1/n\text{ for every }n\in\N_{>0}\}.
\]
Normal forms give
\[
 \Fin=\{v\ge0\},\qquad \mm=\{v>0\},\qquad
 \Fin/\mm\simeq\C.
\]
The standard-part homomorphism is coefficient extraction at exponent
zero:
\[
 \st:\Fin\longrightarrow\C,
 \qquad z=\st(z)+\epsilon,\quad\epsilon\in\mm.
\]
For $a\in\C$ write $\mu(a)=a+\mm$. More generally, $c+\mm$ is a
well-defined infinitesimal neighborhood for any $c\in\SC$, although
$c$ need not have a standard part.

\begin{definition}[Tube over a complex domain]\label{def:tube}
For $U\subseteq\C$ open, its finite tube is
\[
 U^\sharp=\st^{-1}(U)\subseteq\Fin.
\]
For $c\in\SC$ and $r\in\SC^*$, a scale chart is
$c+rU^\sharp$, with normalized coordinate $w=(z-c)/r$.
\end{definition}
For example, $\disk^\sharp$ is \emph{not} the entire surcomplex ball
$\{|z|<1\}$: the latter also contains $1-t$, whose standard part lies
on the ordinary boundary. This distinction matters in Schwarz-type
statements.

\subsection{Set-sized coefficient fields}
For a set-sized divisible ordered subgroup $\Gamma\subseteq\No$,
let
\[
 K_\Gamma=\C((t^\Gamma)),\qquad
 R_\Gamma=\R((t^\Gamma)).
\]
The canonical normal-form embeddings place these fields in $\SC$ and
$\No$, respectively. The standard Hahn-field theorem gives that
$R_\Gamma$ is real closed and $K_\Gamma$ is algebraically closed
\cite{Poonen}. In particular, all roots of a polynomial over
$K_\Gamma$ already belong to $K_\Gamma$.

\begin{proposition}[Containment of set-sized data]\label{prop:fragment}
Every set of surcomplex numbers is contained in some $K_\Gamma$.
\end{proposition}
\begin{proof}
Take the union of the supports of its elements and then the divisible
additive subgroup of $\No$ generated by that union. Both operations
produce sets. Every original normal form has support in this group.
\end{proof}
Consequently the roots, coefficients, and support recursions in any
single construction below can be handled in set-sized algebra. We do
not sum over the whole class of surreal exponents.

\section{Why ordinary topology is not enough}\label{sec:topology}
The \emph{fine topology} on $\SC$ has balls
$B(a,r)=\{z:|z-a|<r\}$ for arbitrary positive $r\in\No$.
Limits and derivatives in this topology are understood directly in
$\epsilon$--$\delta$ form, not by a sequential criterion.

\begin{theorem}[Set-indexed convergence collapses]\label{thm:nets}
Every set-indexed net in $\SC$ converging in the fine topology is
eventually equal to its limit. Every set of surcomplex points is a
closed discrete subspace of $\SC$.
\end{theorem}
\begin{proof}
If $D$ is a set of positive surreal numbers, the surreal cut
$\{0\mid D\}$ gives a positive number smaller than every member of
$D$. For a net $(z_j)_{j\in J}$ and a proposed limit $z$, apply this
to the set of nonzero distances $|z_j-z|$. A sufficiently small ball
about $z$ contains none of the unequal values. Convergence therefore
forces eventual equality.

For a set $A$, the same construction, applied to distances from any
point outside $A$, proves closedness. Applied to the nonzero pairwise
distances within $A$, it gives a common positive separation radius,
and hence discreteness. Empty distance sets cause no difficulty.
\end{proof}
In particular, the subspace topology induced on $\C$ is discrete, not
its Euclidean topology. A nonconstant ordinary complex contour is not
being regarded below as a continuous map from a real interval into
this fine topology.

\begin{proposition}[Clopen infinitesimal neighborhoods]\label{prop:clopen}
Every coset of $\mm$ is both open and closed. The fine topology is
totally disconnected.
\end{proposition}
\begin{proof}
An infinitesimal ball about a point in a coset remains in that coset.
The other cosets are open as well, proving closedness. Given distinct
$z,w$, the scaled coset $z+(z-w)\mm$ is clopen and contains $z$ but
not $w$. Thus no connected subclass containing both points exists.
\end{proof}

\begin{example}[Analytic functions invisible to differentiation]
The function $\chi_{\mm}$, equal to $1$ on $\mm$ and $0$ elsewhere,
is locally constant on all of $\SC$. It is everywhere locally given
by a constant power series, has derivative zero, is bounded by $1$,
and is not constant. On $\Fin$, the map $z\mapsto\st(z)$ is also
locally constant. It agrees with the identity on $\C$, whereas its
derivative is zero and the derivative of $z\mapsto z$ is one.
\end{example}
These examples do not contradict the classical identity theorem:
ordinary connected domains have disappeared. They show that the
fine-topological local definition cannot, by itself, replace the
entire package of classical complex analysis.

\section{Hahn summation and admissible substitutions}\label{sec:sums}
\begin{definition}[Hahn-summable family]\label{def:summable}
A set-indexed family $(z_j)_{j\in J}$ of normal forms is
\emph{Hahn summable} if the union of its supports is well ordered and
each exponent belongs to only finitely many of those supports. Its
Hahn sum $\Hs_j z_j$ is formed coefficientwise.
\end{definition}
No infinite addition of complex coefficients at one exponent is allowed
by this definition. This is why $\sum_{n\ge0}1/n!$, regarded as a
family of constant Hahn series, is not a Hahn sum. Its classical value
$e$ is instead used as a coefficient.

We use the following precise form of the Hahn--Neumann lemma
\cite{Neumann,PoonenUnits}.
\begin{lemma}[Support calculus]\label{lem:neumann}
Let $A,B$ be well-ordered subsets of an ordered abelian group. Then
$A+B$ is well ordered, and each element has only finitely many
representations $a+b$ with $a\in A$, $b\in B$. If $E$ is well
ordered and strictly positive, its additive monoid
\[
 M(E)=\{0\}\cup\bigcup_{n\ge1}(E+\cdots+E)
\]
is well ordered. Every element has only finitely many representations
as a finite ordered list of elements of $E$.
\end{lemma}
Repeated entries are permitted in the lists, and the empty list represents
zero. In particular, the lemma supplies both support well-ordering and
coefficient finiteness; checking just the first is insufficient.

\begin{corollary}[Formal series at infinitesimals]\label{cor:formalcomplex}
Every $P(X)\in\C[[X]]$ can be evaluated by Hahn summation at every
$\epsilon\in\mm$. Finite multivariable series can likewise be
evaluated at infinitesimal tuples. Products, formal substitutions with
zero constant term, and regroupings of jointly summable families have
their expected values.
\end{corollary}
\begin{proof}
Take the union $E$ of the positive supports of the finitely many
increments. Every term has support in $M(E)$, and the finite-list
conclusion of Lemma~\ref{lem:neumann} proves coefficient finiteness.
Formal identities then hold coefficient by coefficient, where each
calculation is finite.
\end{proof}
For example,
\[
 \exp_{\mm}(\epsilon)=\Hs_{n\ge0}\frac{\epsilon^n}{n!},
 \qquad
 \log_{\mm}(1+\epsilon)=\Hs_{n\ge1}
 \frac{(-1)^{n+1}}{n}\epsilon^n
\]
are well defined and inverse to one another.

\begin{example}[A sum without convergence]
With $t=\omega^{-1}$,
\[
 \Hs_{n\ge0}t^n=(1-t)^{-1}.
\]
Its partial sums are not eventually constant and hence do not converge
in the full fine topology. More concretely, their nonzero remainders
have valuation $n+1$, whereas $t^\omega$ is smaller than every such
remainder. Even within a fixed high-rank Hahn field, existence of a
Hahn sum need not imply convergence in its valuation topology.
\end{example}

\begin{lemma}[Evaluation below a coefficient field]\label{lem:ultrasmall}
Let $A_n\in K_\Gamma$ for $n\ge0$. If $h\ne0$ and
\[
 v(h)>\gamma\quad\text{for every }\gamma\in\Gamma,
\]
then $(A_nh^n)_{n\ge0}$ is Hahn summable. Such $h$ fill a nonempty
fine neighborhood of zero, apart from zero itself.
\end{lemma}
\begin{proof}
Write $h=c t^\delta(1+u)$ with $c\in\C^*$, $v(u)>0$, and
$\delta=v(h)>\Gamma$. The set
\[
 B=\bigcup_{n\ge0}\bigl(\supp A_n+n\delta\bigr)
\]
is an ordered concatenation of well-ordered sets: all elements in its
$n$th block precede all elements in its $(n+1)$st block, since
$\delta>\Gamma$. A member of $B$ determines its block uniquely.
Expand $(1+u)^n$ by the finite binomial theorem. The total support is
contained in $B+M(\supp u)$, which is well ordered. For a fixed
exponent there are finitely many decompositions into a member of $B$
and a member of this monoid, finitely many lists from the support of
$u$, and finitely many binomial terms. Thus coefficient finiteness
also holds.

Because $\Gamma$ is a set, choose $\Delta\in\No$ larger than every
member of $\Gamma$. Every nonzero $h$ with $|h|<t^\Delta$ has
$v(h)\ge\Delta>\Gamma$.
\end{proof}
\begin{lemma}[Joint summation after a smaller substitution]\label{lem:jointsmall}
Suppose $(A_{j,k})_j$ is Hahn summable in $K_\Gamma$ for each
$k\ge0$. If $v(h)>\Gamma$, the full family
$(A_{j,k}h^k)_{j,k}$ is jointly Hahn summable.
\end{lemma}
\begin{proof}
For each $k$ use the union of the supports of all $A_{j,k}$, before
any cancellation, as the $k$th support block in the preceding proof.
It is well ordered by the hypothesis, and each exponent in it has
only finitely many contributing $j$. The disjoint translated blocks
and the positive-support binomial expansion of $h^k$ give exactly the
same finite-decomposition argument. Thus regrouping by $k$ loses no
summability information.
\end{proof}
The same arguments apply to finitely many variables whose nonzero
increments all have valuation greater than $\Gamma$. To see the support
control explicitly, choose an increment $h$ of maximal modulus and write
each increment as $h u_j$, where $u_j$ is finite. Group by total degree;
there are finitely many monomials at each degree. The coefficients shifted
by $n v(h)$ form the same ordered support blocks. The positive supports
of the $u_j$ and of the unit part of $h$ generate one monoid, so the
full support lies in the sum of these blocks and that monoid. The
finite-decomposition argument proves joint summability. If one first
makes a summable substitution at $h_0$, enlarge the coefficient field
to contain its support, then use Lemma~\ref{lem:jointsmall} for a
still smaller increment. These are the mechanisms behind the
formal-germ manipulations below.

\section{Local analytic germs}\label{sec:germs}
\subsection{Every formal series has a local realization}
\begin{definition}
A surcomplex analytic germ at $a\in\SC$ is the germ of a function
locally representable as
\[
 F(a+h)=\Hs_{n\ge0} A_n h^n,\qquad A_n\in\SC,
\]
on a fine neighborhood where the family is Hahn summable. Whenever a
formal series is used as a germ, its evaluation neighborhood may be
shrunk to one supplied by Lemma~\ref{lem:ultrasmall}.
\end{definition}

\begin{theorem}[Taylor calculus for germs]\label{thm:germcalculus}
Every formal series in $\SC[[X]]$ realizes an analytic germ. The
realization is injective, respects addition, multiplication, and
composition with zero-constant-term germs, and satisfies
\[
 \frac{d}{dh}\Hs_{n\ge0}A_nh^n
   =\Hs_{n\ge1}nA_nh^{n-1}.
\]
Consequently the ring of these germs at a fixed center is naturally
isomorphic to $\SC[[X]]$, and $A_n=F^{(n)}(a)/n!$.
\end{theorem}
\begin{proof}
Contain the coefficients in one $K_\Gamma$ and use
Lemma~\ref{lem:ultrasmall}. Formal ring and composition identities
involve finite coefficient calculations; simultaneous sufficiently
small substitutions make the expanded identities jointly summable.

Here is the essential estimate behind differentiation. If all
coefficients of $R(X)$ lie in $K_\Gamma$ and $v(h)>\Gamma$, then
\[
 R(h)=R(0)+o_{K_\Gamma}(1),
\]
where $o_{K_\Gamma}(1)$ means smaller in modulus than every positive
element of $R_\Gamma$. Indeed, every term of positive degree has
valuation greater than every member of $\Gamma$, and the same holds
for their Hahn sum. Apply this to
\[
 P(h)-P(0)-A_1h=h^2\sum_{n\ge0} A_{n+2}h^n.
\]
The last factor is bounded, uniformly on such increments, by the fixed
surreal number $|A_2|+1$. The difference quotient therefore tends to
$A_1$ in the full $\epsilon$--$\delta$ sense.

At another sufficiently small center $h_0$, the binomial identity
reexpands $P(h_0+k)$ as a series in $k$. The coefficients are
Hahn sums in $h_0$. Enlarge the coefficient field to contain their
supports and take $k$ smaller still. The same estimate proves the
derivative formula there. Repetition gives all higher derivatives.
For each fixed Taylor degree in $k$, the sum over the original
degree is summable at $h_0$. Lemma~\ref{lem:jointsmall}, in the
enlarged coefficient field, makes the full binomial expansion jointly
summable for the smaller $k$. No limit of partial sums is used.

Finally, a nonzero series with first nonzero term $A_mX^m$ has, on
sufficiently small nonzero increments,
\[
 P(h)=A_m h^m(1+\eta_h),\qquad v(\eta_h)>0.
\]
It cannot represent the zero germ. This proves injectivity and the
Taylor coefficient assertion.
\end{proof}

\begin{remark}
The series $\sum n!X^n$ has classical radius of convergence zero, yet
it defines a surcomplex analytic function on $\mm$. Thus the word
``analytic'' at this local level describes Hahn evaluation of formal
series, not positive classical radius of convergence.
\end{remark}

\subsection{Cauchy--Riemann equations}
Write $F=u+iv$ and $z=x+iy$. A real differential here is an
$\No$-linear map $\No^2\to\No^2$ with an error $o(|h|)$ in the
fine topology.
\begin{proposition}\label{prop:CR}
An analytic germ has a real differential and obeys
\[
 u_x=v_y,\qquad u_y=-v_x.
\]
Conversely, at a point where the real differential exists, these
equalities make that differential complex linear and imply existence
of the surcomplex derivative. The equations alone, without the real
differentiability hypothesis, are not asserted to characterize the
coherent analytic class introduced later.
\end{proposition}
\begin{proof}
The first-order expansion is
$F(z+h)=F(z)+F'(z)h+o(|h|)$. Multiplication by
$\alpha+i\beta$ has real matrix
$\bigl(\begin{smallmatrix}\alpha&-\beta\\\beta&\alpha\end{smallmatrix}\bigr)$,
which gives the equations. Conversely, the equations force a real
Jacobian to have this matrix form. Dividing the remainder by $h$
then proves complex differentiability.
\end{proof}
For analytic functions, equality of mixed derivatives further gives
$u_{xx}+u_{yy}=v_{xx}+v_{yy}=0$. These are differential identities;
no assertion about a global Dirichlet problem is implicit.

\subsection{Inversion and ramification}
\begin{theorem}[Local inverse and ramification]\label{thm:inverse}
If $F(a+h)=b+A_1h+\cdots$ and $A_1\ne0$, then $F$ has an analytic
inverse between sufficiently small fine neighborhoods of $a$ and $b$.
More generally, a nonconstant germ admits a local coordinate $q$ with
$q(0)=0$, $q'(0)\ne0$, and
\begin{equation}\label{eq:ramification}
 F(a+h)-F(a)=q(h)^m,
\end{equation}
where $m$ is its first nonzero Taylor degree after the constant term.
It is open locally and takes each sufficiently small nonzero target
increment exactly $m$ times in a suitable neighborhood.
\end{theorem}
\begin{proof}
The formal inverse coefficients are determined recursively from
$P(Q(Y))=Y$, beginning with $Q'(0)=A_1^{-1}$; the coefficient of each
new unknown is $A_1$. The resulting series and the composition
identities are realizable by Theorem~\ref{thm:germcalculus}.

For clarity, the formal inverse supplies a topological inverse, not
just an identity of symbols. On sufficiently small $h,k$, the
factorization
\[
 P(h)-P(k)=(h-k)(A_1+E(h,k))
\]
has $|E(h,k)|<|A_1|/2$, giving injectivity. Shrink the target ball so
that the evaluated formal inverse lies in this source neighborhood.
The two composition identities then give a bijection onto that ball,
with continuous analytic inverse.

For the general case write
$F(a+h)-F(a)=A_mh^m(1+R(h))$, $R(0)=0$.
Choose $\alpha\in\SC$ with $\alpha^m=A_m$, and put
\[
 q(h)=\alpha h\sum_{n\ge0}{1/m\choose n}R(h)^n.
\]
This is an analytic germ, $q'(0)=\alpha\ne0$, and its $m$th power
is the original series. On a ball where $q$ is invertible, every
sufficiently small nonzero target has exactly $m$ distinct $m$th
roots in the image ball, since $\SC$ is algebraically closed and
has characteristic zero. Pulling them back proves the assertion.
\end{proof}

\section{Canonical lifting of classical holomorphic functions}\label{sec:lift}
\begin{definition}[Taylor lift]
For $f\in\OO(U)$ define
\begin{equation}\label{eq:lift}
 f^\sharp(a+\epsilon)=
 \Hs_{n\ge0}\frac{f^{(n)}(a)}{n!}\epsilon^n,
 \qquad a\in U,\quad\epsilon\in\mm.
\end{equation}
The decomposition into standard part and infinitesimal is unique, so
this is a canonical function on $U^\sharp$.
\end{definition}

\begin{theorem}[Functorial Taylor lifting]\label{thm:lift}
Taylor lifting is an injective $\C$-algebra homomorphism, commutes
with differentiation, and preserves composition:
\[
 (f+g)^\sharp=f^\sharp+g^\sharp,\quad
 (fg)^\sharp=f^\sharp g^\sharp,\quad
 (f')^\sharp=(f^\sharp)',\quad
 (f\circ g)^\sharp=f^\sharp\circ g^\sharp.
\]
The last identity is understood when $g:U\to V$ and $f\in\OO(V)$.
It also preserves reciprocals on zero-free domains and satisfies
$\st(f^\sharp(z))=f(\st z)$.
\end{theorem}
\begin{proof}
All evaluations exist by Corollary~\ref{cor:formalcomplex}.
Multiplication is the Taylor product formula. For composition,
$g^\sharp(a+\epsilon)-g(a)$ is infinitesimal, and substitution of
the two classical Taylor series is jointly Hahn summable. Its formal
coefficients are those of $f\circ g$. Differentiation follows from
Theorem~\ref{thm:germcalculus}; the formal translation formula makes
the same statement at every point of $\mu(a)$, not just at $a$.
Reciprocals follow by applying the product assertion to $f(1/f)=1$.
The standard-part assertion is the constant term of \eqref{eq:lift},
and restriction to ordinary points proves injectivity.
\end{proof}

\begin{proposition}[Zeros and local mapping of a lift]\label{prop:liftzeros}
If $f$ is not identically zero on connected $U$, the zeros of
$f^\sharp$ are exactly the ordinary zeros of $f$, with the same
multiplicities. A classical biholomorphism $f:U\to V$ lifts to a
bijection $U^\sharp\to V^\sharp$ with inverse $(f^{-1})^\sharp$.
\end{proposition}
\begin{proof}
If $a$ is a zero of multiplicity $m$, write
$f(z)=(z-a)^m u(z)$ on a disk about $a$, with $u(a)\ne0$. Then
\[
 f^\sharp(a+\epsilon)=\epsilon^m u^\sharp(a+\epsilon),
\]
and the second factor is a unit with nonzero standard part. Thus only
$\epsilon=0$ is a zero. A monad over a nonzero value has no zero by
standard parts. The inverse assertion follows by lifting the two
composition identities.
\end{proof}
A canonical lift alone therefore creates no infinitesimally displaced
zeros. To study genuine surcomplex perturbations we need surcomplex
coefficients as well.

\section{Coherent Hahn-holomorphic families}\label{sec:H}
Fix a connected ordinary domain $U$ and a set-sized divisible
$\Gamma\subseteq\No$.

\begin{definition}\label{def:H}
An element of $\HH_\Gamma(U)$ is a formal Hahn family
\[
 F=\sum_{\gamma\in S} f_\gamma t^\gamma,
 \qquad f_\gamma\in\OO(U),
\]
with $S=\{\gamma:f_\gamma\not\equiv0\}$ well ordered. Its evaluation
at $z=a+\epsilon\in U^\sharp$ is
\begin{equation}\label{eq:Heval}
 F(z)=\Hs_{\gamma\in S}\Hs_{n\ge0}
 t^\gamma\frac{f_\gamma^{(n)}(a)}{n!}\epsilon^n.
\end{equation}
We suppress a separate evaluation symbol once this meaning is fixed.
For $F\ne0$, define its family valuation and leading function by
\[
 v_H(F)=\min S,\qquad \LC_H(F)=f_{\min S}.
\]
These are not the same as the pointwise valuation and leading
coefficient at a zero of the leading function.
\end{definition}
The entire family of coefficients, including their support, is part of
the analytic structure. The functions $f_\gamma$ need not be uniformly
bounded as $\gamma$ varies. No such condition is required for Hahn
summability.

\begin{theorem}[Evaluation, Taylor expansion, and algebra]\label{thm:Heval}
Equation~\eqref{eq:Heval} is jointly Hahn summable, defines an analytic
function in the fine local sense, and is injective as a representation
by functions. Evaluation commutes with addition, multiplication, and
coefficientwise differentiation. If $z_0,z\in\mu(a)$, then
\begin{equation}\label{eq:monadTaylor}
 F(z)=\Hs_{n\ge0}\frac{F^{(n)}(z_0)}{n!}(z-z_0)^n.
\end{equation}
This exact expansion is valid throughout that monad, not just on a
smaller ball.
\end{theorem}
\begin{proof}
Let $E=\supp\epsilon\subseteq\No_{>0}$. The support of the double
family is contained in $S+M(E)$. Lemma~\ref{lem:neumann} gives
well-ordering, finitely many decompositions $\lambda=\gamma+\mu$,
and finitely many lists representing $\mu$. These give exactly the
two requirements for joint summability.

For \eqref{eq:monadTaylor}, put $z_0=a+\epsilon_0$ and $z=z_0+h$.
Expand each classical Taylor series at $a$ in the two infinitesimal
variables $\epsilon_0,h$. The entire expansion has support in
$S+M(\supp\epsilon_0\cup\supp h)$ and is jointly summable. The
binomial regrouping is therefore legitimate and gives the stated
coefficients.

To see directly that the derivative is a genuine fine-topological
limit, let $\gamma_0=\min S$. The Taylor remainder can be written
\[
 F(z_0+h)-F(z_0)-hF'(z_0)=h^2t^{\gamma_0}R(\epsilon_0,h).
\]
For all infinitesimal $h$, $R$ is finite and its standard part is the
fixed complex number $f_{\gamma_0}''(a)/2$. Hence
$|R|<|f_{\gamma_0}''(a)|/2+1$. Given any positive surreal error bound,
choose $|h|$ small enough to bound
$|h|t^{\gamma_0}|R|$ by that error. This proves differentiation;
higher derivatives work identically. The same support argument proves
products.

Finally, for an ordinary $a\in U$, the normal form of $F(a)$ is
$\sum f_\gamma(a)t^\gamma$. If all evaluations vanish, every
$f_\gamma(a)$ vanishes for all $a$, so the family is zero.
\end{proof}

\begin{proposition}[Admissible composition]\label{prop:Hcompose}
Let $F\in\HH_\Gamma(V)$ and
\[
 G=g_0+E\in\HH_\Gamma(U),\qquad
 g_0:U\to V\text{ holomorphic},\qquad v_H(E)>0.
\]
Then $F\circ G$ belongs to $\HH_\Gamma(U)$, possibly after enlarging
$\Gamma$ to include all the stated coefficients, and is represented by
\begin{equation}\label{eq:Hcompose}
 \sum_{\gamma\in\supp F}\sum_{n\ge0}t^\gamma
 \frac{f_\gamma^{(n)}\circ g_0}{n!}E^n.
\end{equation}
Its evaluation is the composition of the evaluated functions.
\end{proposition}
\begin{proof}
The positive support of $E$ generates a well-ordered monoid, and the
support of \eqref{eq:Hcompose} is contained in its sum with $\supp F$.
The same finite-decomposition lemma proves coefficient finiteness.
The coefficients are holomorphic compositions and finite sums of
products. At a finite point, $G$ has standard part $g_0(\st z)$;
the formal Taylor identity and joint summability identify all
coefficients of the two evaluated expressions.
\end{proof}

\begin{proposition}[Units and the leading function]\label{prop:unitforward}
If $F=t^{\gamma_0}(f_0+E)$, $v_H(E)>0$, and $f_0$ is zero free
on $U$, then $F$ is a unit of $\HH_\Gamma(U)$ and never vanishes
on $U^\sharp$. Its inverse is
\[
 F^{-1}=t^{-\gamma_0}f_0^{-1}
       \sum_{n\ge0}(-E/f_0)^n.
\]
\end{proposition}
\begin{proof}
The support of $E/f_0$ is positive. The geometric series is a
well-defined Hahn family by Lemma~\ref{lem:neumann}, with holomorphic
coefficients. Its product with $F$ is one. Pointwise, the normalized
factor has standard part $f_0(\st z)\ne0$.
\end{proof}

\subsection{What coherence supplies}
The family $f_\gamma$ is a single holomorphic function on the whole
connected base, not a separately chosen germ at each $a$. Thus
$\st(z)$ on $U^\sharp$ is not in $\HH_\Gamma(U)$ when $U$ has more
than one point. If it were, its values at ordinary points would force
it to be the canonical lift of the identity, which it is not.

On connected domains restriction maps are injective. Coefficientwise
families on overlapping connected patches that agree as evaluated
functions have equal coefficients on overlaps, and these coefficients
glue holomorphically. Their nonzero support sets coincide along any
chain of overlapping connected patches, by the ordinary identity
theorem. A connected cover therefore glues to a family of the same
well-ordered support. On a disconnected base one should impose this
condition componentwise, rather than silently require a common
well-ordered support for arbitrarily many components.

\section{Identity, open mapping, and maximum modulus}\label{sec:rigidity}
\begin{theorem}[Identity theorem]\label{thm:identity}
Let $U$ be connected and $F\in\HH_\Gamma(U)$. Each of the following
conditions implies $F=0$:
\begin{enumerate}[label=(\roman*)]
\item $F$ vanishes at ordinary points with an ordinary accumulation
point in $U$;
\item all derivatives $F^{(n)}(z_0)$ vanish at one $z_0\in U^\sharp$;
\item $F$ vanishes on a nonempty fine-open subset of $U^\sharp$.
\end{enumerate}
\end{theorem}
\begin{proof}
For (i), uniqueness of normal forms makes every $f_\gamma$ vanish on
the same ordinary set. Apply the classical identity theorem to each
coefficient.

For (ii), let $a=\st z_0$. Equation~\eqref{eq:monadTaylor} first
makes $F$ identically zero on $\mu(a)$. All derivatives at $a$ then
vanish, so uniqueness of the normal forms of $F^{(n)}(a)$ gives
$f_\gamma^{(n)}(a)=0$ for every $\gamma,n$. Each coefficient has
zero classical Taylor series at $a$ and hence is zero on connected
$U$. Condition (iii) gives (ii) at an interior point by repeated
differentiation of the zero function.
\end{proof}

\begin{corollary}[Open mapping and isolated zeros]\label{cor:open}
A nonconstant element of $\HH_\Gamma(U)$ is an open map
$U^\sharp\to\SC$ for the fine topologies. A nonzero element has
isolated zeros in that topology. If $F'=0$, then $F$ is a constant
surcomplex number.
\end{corollary}
\begin{proof}
At any point, a nonconstant family has a nonconstant germ; otherwise
it would agree with a constant on a fine-open set, and the identity
theorem would make it constant globally. Apply the local ramification
theorem. The same normal form gives isolated zeros. If $F'=0$, every
classical coefficient has derivative zero and is constant on $U$.
\end{proof}
Here openness is genuine even though the source is totally
disconnected. It is a local algebraic phenomenon, not an appeal to
connectedness of the fine topology.

\begin{theorem}[Maximum-modulus principle]\label{thm:maxmod}
If $|F|$ has a local maximum in the fine topology at a point of
$U^\sharp$, where $U$ is connected and $F\in\HH_\Gamma(U)$,
then $F$ is constant.
\end{theorem}
\begin{proof}
Suppose the germ at $z_0$ is nonconstant. Write
\[
 F(z_0+h)=A_0+A_mh^m+\cdots,\qquad A_m\ne0.
\]
If $A_0=0$, sufficiently small nonzero $h$ gives nonzero values,
contradicting maximality. Otherwise choose $d\in\SC^*$ with
$A_md^m=A_0$. With a positive surreal $s$ sufficiently small relative
to all the coefficient data, the local expansion becomes
\[
 F(z_0+ds)=A_0(1+s^m+E_s),\qquad |E_s|<\tfrac12s^m.
\]
The real part of the last parenthesis is greater than one, so its
modulus is greater than one. These increments can be made smaller
than any prescribed fine radius, again contradicting maximality.
Thus the germ is constant. The identity theorem makes the whole
family constant.
\end{proof}
This theorem does not assert that a maximum is attained on a closed
surcomplex disk. The missing compactness is precisely why the
maximum principle need not imply a naive Liouville theorem.

\section{Support-controlled Weierstrass preparation}\label{sec:weierstrass}
This section gives a constructive factorization over the Hahn
coefficients. It is the main algebraic input to zero counting and the
residue theorem.

\subsection{Preparation}
\begin{theorem}[Hahn--Weierstrass preparation]\label{thm:Wprep}
Let $D\subseteq\C$ be a disk about zero and suppose
\begin{equation}\label{eq:Winput}
 F(x)=x^m u_0(x)+\sum_{\gamma\in S} f_\gamma(x)t^\gamma,
 \qquad m\ge1,
\end{equation}
where $u_0,f_\gamma\in\OO(D)$, $u_0$ is zero free on $D$, and
$S\subseteq\Gamma_{>0}$ is well ordered. There exist unique
\[
 W(x)=x^m+b_{m-1}x^{m-1}+\cdots+b_0,
 \qquad b_j\in K_\Gamma,\quad v(b_j)>0\text{ or }b_j=0,
\]
and
\[
 U(x)=u_0(x)+\sum_{\lambda>0}u_\lambda(x)t^\lambda
       \in\HH_\Gamma(D)
\]
such that $F=UW$. The positive supports of $U$ and the $b_j$ are
contained in $M(S)\setminus\{0\}$. The factor $U$ is a unit on
$D^\sharp$.
\end{theorem}
\begin{proof}
Put $M=M(S)$ and extend $f_\lambda$ by zero for
$\lambda\notin S$. Construct
\[
 U=\sum_{\lambda\in M}U_\lambda t^\lambda,
 \qquad W=\sum_{\lambda\in M}W_\lambda t^\lambda,
 \qquad U_0=u_0,\quad W_0=x^m,
\]
where $W_\lambda$ is a polynomial of degree less than $m$ for
$\lambda>0$. Suppose all coefficients of smaller exponent have been
constructed. The coefficient equation at $\lambda>0$ is
\begin{equation}\label{eq:Wlinear}
 x^mU_\lambda+u_0W_\lambda=H_\lambda,
 \qquad
 H_\lambda=f_\lambda-
 \sum_{\substack{\alpha+\beta=\lambda\\\alpha,\beta>0}}
 U_\alpha W_\beta.
\end{equation}
The sum is finite, and all its indices precede $\lambda$.

For any holomorphic $h$, let $\remm h$ be its Taylor polynomial at
zero of degree less than $m$. Set
\begin{equation}\label{eq:Wrecursion}
 W_\lambda=\remm(H_\lambda/u_0),\qquad
 U_\lambda=u_0\frac{H_\lambda/u_0-W_\lambda}{x^m}.
\end{equation}
The apparent singularity at zero is removable, so the new
coefficients are holomorphic throughout $D$. Transfinite recursion
along the well-ordered set $M$ now defines all coefficients.
Lemma~\ref{lem:neumann} ensures finite products at every exponent,
so the constructed product equals $F$ as a Hahn family. Its evaluated
identity follows from Theorem~\ref{thm:Heval}.

For uniqueness, compare two normalized factorizations. The union of
their positive supports and $S$ is still a well-ordered positive set;
work in the monoid it generates. At the first exponent of disagreement,
\eqref{eq:Wlinear} has zero right-hand side. Division by $u_0$ and
separation of the Taylor terms below degree $m$ force both differences
to vanish, a contradiction. Finally $U$ has the zero-free leading
function $u_0$, so Proposition~\ref{prop:unitforward} gives its inverse.
\end{proof}

\begin{remark}[What does not have to converge]
The corrections in \eqref{eq:Wrecursion} need not satisfy analytic
norm estimates as the exponent varies. There may be infinitely many
exponents below a given larger exponent. What matters is that each
coefficient of a product has finitely many contributors, and that all
supports lie in a well-ordered monoid. The proof is not an iterative
Banach-space convergence argument.
\end{remark}

\subsection{Division}
\begin{theorem}[Hahn--Weierstrass division]\label{thm:Wdivision}
Let $W$ be a monic polynomial as in Theorem~\ref{thm:Wprep}. For
any $A\in\HH_\Gamma(D)$ there exist unique
\[
 Q\in\HH_\Gamma(D),\qquad R\in K_\Gamma[x],\quad\deg R<m,
\]
such that $A=QW+R$.
\end{theorem}
\begin{proof}
Multiply $A$ by a monomial to make its support nonnegative. Take the
monoid generated by its positive support and the positive supports
of the coefficients of $W$. At exponent zero, divide the holomorphic
coefficient of $A$ by $x^m$, using its Taylor remainder below degree
$m$. At a later exponent $\lambda$, subtract the already determined
finite sum involving positive coefficients of $W$ and lower
coefficients of $Q$. Divide the resulting holomorphic function
$H_\lambda$ by $x^m$:
\[
 R_\lambda=\remm H_\lambda,\qquad
 Q_\lambda=(H_\lambda-R_\lambda)/x^m.
\]
The quotient is holomorphic at zero. Well-ordered recursion gives
$Q,R$, with support in the chosen monoid. Undo the initial monomial
multiplication. A first-exponent comparison proves uniqueness just as
in the preparation theorem.
\end{proof}

\subsection{Why all prepared roots are infinitesimal}
\begin{lemma}\label{lem:smallroots}
Every root of
$W(x)=x^m+\sum_{j<m}b_jx^j$ with $v(b_j)>0$ is infinitesimal.
There are exactly $m$ roots counted with multiplicity, all in
$K_\Gamma$.
\end{lemma}
\begin{proof}
Algebraic closedness of $K_\Gamma$ supplies all $m$ roots. If a root
$r$ had $v(r)<0$, the leading term $r^m$ would have strictly smaller
valuation than every $b_jr^j$, since
\[
 v(b_jr^j)-v(r^m)=v(b_j)+(j-m)v(r)>0.
\]
It could not be canceled. If $v(r)=0$, reduction modulo the maximal
ideal would give $\st(r)^m=0$, also impossible. Thus $v(r)>0$,
including the possibility $r=0$. Since $K_\Gamma$ already splits
$W$, extension to the full class field introduces no additional roots.
\end{proof}

\section{Specialization of zeros and Rouch\'e comparison}\label{sec:zeros}
\begin{theorem}[Exact specialization of the zero divisor]\label{thm:zeros}
Let $0\ne F\in\HH_\Gamma(U)$ with connected $U$, and write
\[
 F=t^{\gamma_0}(f_0+E),\qquad v_H(E)>0,
\]
where $f_0$ is not identically zero. Then:
\begin{enumerate}[label=(\roman*)]
\item every zero of $F$ lies over a classical zero of $f_0$;
\item for each $a\in U$ with $\ord_a(f_0)=m$, the monad $\mu(a)$
contains exactly $m$ zeros of $F$, counted with analytic multiplicity;
\item these zeros belong to $K_\Gamma$, and the whole zero set is
at most countable.
\end{enumerate}
Equivalently, if zero divisors retain multiplicities, then
\begin{equation}\label{eq:divspecialization}
 \st_*\operatorname{div}_0(F)=\operatorname{div}_0(f_0).
\end{equation}
\end{theorem}
\begin{proof}
At $z=a+\epsilon$, reduction gives
$\st(t^{-\gamma_0}F(z))=f_0(a)$, proving (i).
For a zero $a$ of order $m$, choose a disk about $a$ on which
$f_0(a+x)=x^m u_0(x)$ with $u_0$ zero free. Apply preparation to
$t^{-\gamma_0}F(a+x)$. The unit never vanishes, and the prepared
polynomial has exactly $m$ infinitesimal roots by
Lemma~\ref{lem:smallroots}. Multiplication by the unit preserves
the least nonzero Taylor degree at each root, so polynomial and
analytic multiplicities agree. This proves (ii) and the field
assertion of (iii).

The zeros of a nonzero classical holomorphic function in a planar
domain form a discrete, at most countable set. Each standard-part
fiber contains only finitely many distinct roots. Their union is
therefore at most countable.
\end{proof}

\begin{corollary}[Characterization of units]\label{cor:units}
For $0\ne F\in\HH_\Gamma(U)$, the following are equivalent:
$F$ is a unit of $\HH_\Gamma(U)$; its leading function is zero free
on $U$; and its evaluation has no zero in $U^\sharp$.
\end{corollary}
\begin{proof}
The forward leading-coefficient comparison in $FG=1$ makes the
leading function a unit in $\OO(U)$. Proposition~\ref{prop:unitforward}
gives the converse. The remaining equivalence is
Theorem~\ref{thm:zeros}.
\end{proof}

\begin{theorem}[Infinitesimal Rouch\'e theorem]\label{thm:rouche}
Let $D$ be a bounded Jordan domain whose closure lies in $U$.
Suppose $F\ne0$, its leading function $f_0$ is nonzero on
$\partial D$, and $v_H(G)>v_H(F)$. Then $F$ and $F+G$ have the
same finite number of zeros in $D^\sharp$, counted with
multiplicity. More strongly, their zero multiplicities agree after
pushforward by standard part, fiber by fiber.
\end{theorem}
\begin{proof}
The leading function is unchanged. Apply
Theorem~\ref{thm:zeros} to each zero of $f_0$ in $D$. There are
finitely many because $f_0$ is holomorphic near the compact closure
and has no boundary zeros.
\end{proof}

\begin{corollary}[Comparison at the same leading scale]
If $v_H(F)=v_H(G)$ and their leading functions satisfy
$|g_0(\zeta)|<|f_0(\zeta)|$ on $\partial D$, then $F+G$ and $F$
have equal total zero counts in $D^\sharp$.
\end{corollary}
\begin{proof}
The leading function of $F+G$ is $f_0+g_0$, which is nonzero on the
boundary. Classical Rouch\'e comparison gives equal zero counts for
$f_0+g_0$ and $f_0$. One way to see that comparison is to vary
$f_0+s g_0$ for $0\le s\le1$: it never vanishes on the boundary,
so its integer winding number is constant. The specialization theorem
now gives the surcomplex count.
\end{proof}
The first version preserves individual standard-part fibers; the
second generally preserves only the total count. Neither counts zeros
with infinite standard-scale coordinate. Such zeros require a different
scale chart.

\section{Contours, actual poles, and the residue theorem}\label{sec:residues}
\subsection{Definition of the integral}
Let $\gamma$ be an ordinary piecewise $C^1$ contour in $U$. For
$F=\sum f_\lambda t^\lambda\in\HH_\Gamma(U)$ define
\begin{equation}\label{eq:Hint}
 \HI_\gamma F(z)\,dz
   :=\sum_\lambda t^\lambda\int_\gamma f_\lambda(z)\,dz.
\end{equation}
The right side is a Hahn series supported in $\supp F$. Whenever a Jordan boundary is integrated below, it is assumed to be
piecewise $C^1$. This is a coefficientwise contour integral. It is not an order-topological
Riemann integral, and the ordinary parameter interval is not being
identified with a fine-topological interval.

The integral is linear over $K_\Gamma$, respects concatenation and
reversal, and satisfies integration by parts whenever both sides are
defined. For example, multiplication by a fixed Hahn scalar can be
moved through it because each product coefficient is a finite sum.
Classical Cauchy's theorem applied to every coefficient immediately
gives
\[
 \Hoint_{\partial D}F(z)\,dz=0
\]
when $F$ is holomorphic on a neighborhood of $\overline D$.

\subsection{Meromorphic fractions and local residues}
Because $U$ is connected, $\HH_\Gamma(U)$ is an integral domain:
the leading coefficient of a product is the product of two nonzero
holomorphic functions. Its fraction field embeds in
\[
 \mathcal M(U)((t^\Gamma)),
\]
where $\mathcal M(U)$ is the classical meromorphic function field.
For a fraction $A/B$, expand $B^{-1}$ by its leading function on
ordinary regions where that function is nonzero. The coefficients of
the resulting Hahn family are meromorphic, with possible poles only
at zeros of the leading function of $B$.

We define \eqref{eq:Hint} for such a fraction along any ordinary
contour avoiding those leading zeros. This definition is independent
of the fractional representation: the embedding in the meromorphic
Hahn field is injective, and normal-form coefficients are unique.

At an actual surcomplex point $b$, the local germs of $A$ and $B$
have formal Taylor expansions. Unless the fraction is identically
undefined, their quotient is a formal Laurent series with finite
principal part:
\[
 \frac{A(b+h)}{B(b+h)}=\sum_{n\ge-N} C_nh^n.
\]
Its \emph{actual local residue} is $C_{-1}$. In particular, this
residue is centered at $b$, not at $\st b$.

\begin{example}[A moving pole]
For $R(z)=1/(z-t)$, the actual pole is at $t$. On ordinary
$z\ne0$ its Hahn expansion is
\[
 R(z)=\sum_{n\ge0}\frac{t^n}{z^{n+1}}.
\]
The coefficient functions have poles at the ordinary point zero,
although the actual rational function is regular there. This expansion
must not be used without qualification at $z$ comparable to $t$.
The fraction itself, or its local Laurent series at the relevant
surcomplex point, is the appropriate evaluation there.
\end{example}

\begin{theorem}[Residue theorem for actual surcomplex poles]\label{thm:residue}
Let $A,B\in\HH_\Gamma(U)$, $B\ne0$, and let $D$ be a bounded
Jordan domain with $\overline D\subset U$. Suppose the leading
function of $B$ is nonzero on $\partial D$. Then $A/B$ has only
finitely many poles in $D^\sharp$, and
\begin{equation}\label{eq:residue}
 \Hoint_{\partial D}\frac{A(z)}{B(z)}\,dz
  =2\pi i\sum_{b\in D^\sharp}
          \Res_{z=b}\frac{A(z)}{B(z)}\,dz.
\end{equation}
The sum on the right is an ordinary finite sum over the actual poles.
\end{theorem}
\begin{proof}
There are finitely many classical zeros of the leading function of
$B$ in $D$. By Theorem~\ref{thm:zeros}, each has a finite fiber of
actual zeros, and every pole of the fraction lies in one of these
fibers.

Choose disjoint ordinary small circles around the classical zeros.
Each meromorphic Hahn coefficient of $A/B$ is holomorphic between
the original boundary and these circles. Coefficientwise classical
Cauchy's theorem reduces the left side of \eqref{eq:residue} to the
sum of the circle integrals. It remains to identify one circle integral
with the residues in its monad.

Translate that classical center to zero. Preparation gives
\[
 B=t^\beta UW,
\]
where $U$ is a Hahn-holomorphic unit and $W$ is a monic polynomial
whose roots $r_j$ are infinitesimal. Divide $A/U$ by $W$ using
Theorem~\ref{thm:Wdivision}:
\[
 \frac AB=t^{-\beta}\left(Q+\frac RW\right),
 \qquad Q\in\HH_\Gamma(D_0),\quad R\in K_\Gamma[x],\quad\deg R<\deg W.
\]
The circle integral of $Q$ is zero. Over the algebraically closed
field $K_\Gamma$, finite partial fractions give
\[
 \frac{R(x)}{W(x)}=\sum_j\sum_{k=1}^{m_j}
                    \frac{C_{jk}}{(x-r_j)^k}.
\]
On an ordinary circle about zero, expand
\begin{equation}\label{eq:movingpole}
 \frac1{(x-r_j)^k}
  =\sum_{n\ge0}{k+n-1\choose n}\frac{r_j^n}{x^{k+n}}.
\end{equation}
This is jointly Hahn summable, since $r_j$ is infinitesimal; multiplying
by a fixed $C_{jk}$ preserves summability. The only term in
\eqref{eq:movingpole} with nonzero contour integral is $k=1,n=0$.
Thus the circle integral is
$2\pi i\,t^{-\beta}\sum_j C_{j1}$. But $t^{-\beta}C_{j1}$ is
exactly the actual local residue of $A/B$ at $r_j$: $Q$ is analytic
there and the other partial fractions are regular or have the
indicated principal parts. This proves the local identification and
hence the theorem.
\end{proof}
The preparation and division steps are what make this more than a
renaming of the classical coefficientwise residue theorem. They
identify coefficientwise contour data with poles at displaced
surcomplex locations.

\subsection{Cauchy's formula and the argument principle}
\begin{corollary}[Cauchy integral formula]\label{cor:cauchy}
For $F\in\HH_\Gamma(U)$, $\overline D\subset U$ as above,
$z\in D^\sharp$, and $n\ge0$,
\begin{equation}\label{eq:cauchy}
 F^{(n)}(z)=\frac{n!}{2\pi i}
 \Hoint_{\partial D}\frac{F(\zeta)}{(\zeta-z)^{n+1}}\,d\zeta.
\end{equation}
\end{corollary}
\begin{proof}
Enlarge $\Gamma$ to include the coefficients of the fixed point $z$.
The denominator belongs to the Hahn family ring and its leading
function is $(\zeta-\st z)^{n+1}$, which has no boundary zero.
The only possible actual pole is at $\zeta=z$, and its residue is
$F^{(n)}(z)/n!$ by the local Taylor expansion, with residue zero
when the singularity is removable. Apply
Theorem~\ref{thm:residue}.
\end{proof}
For $z=a+\epsilon$, the same formula can be checked directly by the
jointly summable kernel expansion
\[
 (\zeta-z)^{-1}
   =\sum_{n\ge0}\frac{\epsilon^n}{(\zeta-a)^{n+1}}.
\]
The coefficient integrals give the Taylor coefficients at $a$, exactly
as required by the definition of evaluation.

\begin{corollary}[Argument principle]\label{cor:argument}
If $F\in\HH_\Gamma(U)$ has leading function nonzero on
$\partial D$, then
\[
 \frac1{2\pi i}\Hoint_{\partial D}\frac{F'(z)}{F(z)}\,dz
    =N_{D^\sharp}(F)\in\N,
\]
where zeros are counted with multiplicity. For a nonzero meromorphic
fraction $R=A/B$ whose numerator and denominator leading functions
are nonzero on the boundary, the corresponding integral of $R'/R$
is the number of zeros minus the number of poles, after cancellation.
\end{corollary}
\begin{proof}
Locally write $F(b+h)=h^mV(h)$ with $V(0)\ne0$. Then
\[
 \frac{F'}F=\frac mh+\frac{V'}V,
\]
so the actual residue is $m$. For a pole, the same computation gives
a negative multiplicity. The residue theorem completes the proof.
\end{proof}
An integer survives even when individual zero positions and individual
residues involve several non-Archimedean levels of scale.

\subsection{Infinitesimal contour deformations}
A lifted parametrization has the form
\[
 \widetilde\gamma(s)=\gamma_0(s)+\eta(s),\qquad
 \eta(s)=\sum_{\lambda\in E}\eta_\lambda(s)t^\lambda,
 \qquad E\subseteq\Gamma_{>0}\text{ well ordered},
\]
where $\gamma_0$ and all $\eta_\lambda$ are piecewise $C^1$ and
have matching endpoints for a closed contour. Define its integral by
Taylor substitution and the pullback
$R(\widetilde\gamma(s))\widetilde\gamma'(s)\,ds$, followed by
coefficientwise integration. The common positive support is essential.

\begin{proposition}[Deformation invariance]\label{prop:deformation}
Suppose $R$ is a meromorphic Hahn fraction and its denominator's
leading function has no zero on $\gamma_0$. Then
\[
 \Hoint_{\widetilde\gamma}R(z)\,dz
      =\Hoint_{\gamma_0}R(z)\,dz.
\]
\end{proposition}
\begin{proof}
Use the homotopy $\gamma_u=\gamma_0+u\eta$ for ordinary
$0\le u\le1$. Support calculus makes every substitution admissible.
At each exponent the expression is a finite polynomial in $u$ with
piecewise continuous coefficients in $s$. The chain rule gives
\[
 \frac{\partial}{\partial u}
  \bigl(R(\gamma_u)\gamma_u'\bigr)
     =\frac{d}{ds}\bigl(R(\gamma_u)\eta\bigr).
\]
Integrate coefficientwise in $s$. The endpoint terms vanish because
the parametrizations are closed, so the integral is independent of
$u$. Evaluating at zero and one proves the assertion.
\end{proof}
This is invariance under a specified algebraic-analytic deformation of
a classical contour. It is not an unrestricted homotopy theorem for
all maps into the fine-topological class field.

\section{Cauchy estimates, Liouville, and Schwarz--Pick}\label{sec:bounds}
\subsection{Coefficient majorants}
For compact $K\subset U$ define
\begin{equation}\label{eq:majorant}
 \mathcal M_K(F)=\sum_{\gamma\in S}
          \left(\sup_{a\in K}|f_\gamma(a)|\right)t^\gamma
          \in R_\Gamma.
\end{equation}
For real Hahn series, $A\cw B$ means that every coefficient of
$B-A$ is nonnegative. This is a partial order much stronger than the
ordinary surreal order. The majorant is not the pointwise modulus of
$F$ and is not claimed to be a norm over $K_\Gamma$.

\begin{proposition}[Coefficientwise Cauchy estimates]\label{prop:cauchybound}
If $\overline{B(a,r)}\subset U$ with $a\in\C$ and ordinary
$r>0$, then
\[
 \mathcal M_{\{a\}}(F^{(n)})
     \cw \frac{n!}{r^n}\mathcal M_{\{|\zeta-a|=r\}}(F).
\]
Moreover,
\[
 \mathcal M_K(F+G)\cw\mathcal M_K(F)+\mathcal M_K(G),\qquad
 \mathcal M_K(FG)\cw\mathcal M_K(F)\mathcal M_K(G).
\]
\end{proposition}
\begin{proof}
Apply the ordinary complex Cauchy estimate separately to each
$f_\gamma$. The first inequality is exactly the collection of those
inequalities. The sum bound is the triangle inequality coefficient by
coefficient. Each product coefficient is a finite convolution, so the
same triangle inequality and the bound for products give the final
assertion.
\end{proof}

\begin{theorem}[Coefficientwise Liouville theorem]\label{thm:liouville}
Let $F\in\HH_\Gamma(\C)$. If each coefficient function is bounded
on $\C$ by some finite real constant, then $F$ is a constant element
of $K_\Gamma$.
\end{theorem}
\begin{proof}
For each exponent $\gamma$, Cauchy's estimate gives
$|f_\gamma'(a)|\le M_\gamma/r$ for every ordinary $r>0$.
Letting $r$ increase within $\R$ forces $f_\gamma'(a)=0$.
Every coefficient is therefore constant. No uniform bound in
$\gamma$ is required.
\end{proof}

\begin{example}[Why a single bound fails]\label{ex:liouvillefail}
The nonconstant function $F(z)=tz$ belongs to $\HH_\Gamma(\C)$
and satisfies $|F(z)|<1$ for every $z\in\C^\sharp=\Fin$.
Indeed, a finite $z$ is bounded by an integer, whereas $t$ is smaller
than the reciprocal of every integer. Nevertheless $F'=t\ne0$.
A single bound hides lower-order coefficients. The classical argument
``smaller than $M/r$ for all real $r$ means zero'' is false for a
surreal derivative: it may only mean infinitesimal.

This is a failure on the finite tube. Separately,
$\chi_{\mm}$ from Section~\ref{sec:topology} is a bounded,
nonconstant locally analytic function on the \emph{whole} class
$\SC$. The two examples concern different analytic classes and
different domains.
\end{example}

\subsection{Exact disk geometry for canonical lifts}
\begin{theorem}[Lifted Schwarz--Pick]\label{thm:SP}
Let $f:\disk\to\disk$ be a classical holomorphic map. Then for
all $z,w\in\disk^\sharp$,
\begin{equation}\label{eq:SP}
 \abs{\frac{f^\sharp(z)-f^\sharp(w)}
 {1-\overline{f^\sharp(w)}f^\sharp(z)}}
 \le
 \abs{\frac{z-w}{1-\overline w z}}.
\end{equation}
For distinct $z,w$, equality holds if and only if $f$ is a classical
disk automorphism. In particular, if $f(0)=0$, then
\[
 |f^\sharp(z)|\le|z|,\qquad |(f^\sharp)'(0)|\le1,
\]
with the usual rotation equality case.
\end{theorem}
\begin{proof}
All denominators in \eqref{eq:SP} are units, because their standard
parts are nonzero. For $z\ne w$, write the ratio of the two
complex expressions before taking moduli as
\[
 Q(z,w)=\frac{f^\sharp(z)-f^\sharp(w)}{z-w}
          \frac{1-\overline w z}
          {1-\overline{f^\sharp(w)}f^\sharp(z)}.
\]
The ordinary divided difference
$(f(\zeta)-f(\xi))/(\zeta-\xi)$ extends holomorphically across
the diagonal. Its two-variable Taylor lift equals the first factor,
by the formal divided-difference identity. Thus $Q$ is finite, and
its standard part is determined even if $z-w$ is infinitesimal.

Put $a=\st z$, $b=\st w$. If $a\ne b$, ordinary Schwarz--Pick
says $|\st Q|<1$ unless $f$ is an automorphism. If $a=b$, the
standard part is
\[
 f'(a)\frac{1-|a|^2}{1-|f(a)|^2},
\]
and the infinitesimal form of classical Schwarz--Pick gives the same
strict inequality unless $f$ is an automorphism. A finite number with
standard modulus strictly less than one itself has modulus less than
one. Multiplying by the nonzero right-hand expression proves strict
inequality in both cases.

For an automorphism, the equality follows from the algebraic identity
for disk M\"obius transformations, which remains valid over the real
closed field $\No$. The case $z=w$ is immediate. Taking $w=0$
gives Schwarz's lemma; the derivative assertion and rotation case
follow from the classical coefficient at zero.
\end{proof}

\begin{example}[Why arbitrary deformations do not satisfy Schwarz]
The Hahn family $F(z)=(1+t)z$ maps $\disk^\sharp$ into itself and
fixes zero, but $|F'(0)|=1+t>1$. Its reduction is the identity disk
map; the infinitesimal dilation is undetectable by standard parts.
It does not map the entire ball $|z|<1$ into itself. Thus the
hypotheses of Theorem~\ref{thm:SP}, particularly canonical lifting,
are substantive.
\end{example}

\section{Primitives, Morera, and differential equations}\label{sec:primitive}
\subsection{Periods and logarithmic primitives}
\begin{theorem}[Primitive criterion]\label{thm:primitive}
For $F\in\HH_\Gamma(U)$ on connected $U$, the following are
equivalent: $F=G'$ for some $G\in\HH_\Gamma(U)$; and
$\Hoint_\gamma F(z)\,dz=0$ for every ordinary closed contour
$\gamma$ in $U$. In particular, every such $F$ has a primitive
when $U$ is simply connected. A prescribed value in $K_\Gamma$ at one ordinary
point makes the primitive unique.
\end{theorem}
\begin{proof}
A derivative has zero periods by coefficientwise integration. In the
other direction, normal-form uniqueness makes every coefficient
period zero. Fix $a_0\in U$ and let
\[
 g_\gamma(a)=\int_{a_0}^a f_\gamma(\zeta)\,d\zeta,
\]
using ordinary path independence. Then $g_\gamma$ is holomorphic,
$g_\gamma'=f_\gamma$, and $g_\gamma(a_0)=0$. The family
$G=\sum g_\gamma t^\gamma$ has support contained in that of $F$
and is a primitive. A prescribed constant can be added. Uniqueness
follows because a family with zero derivative is constant.
\end{proof}

If $F=t^\beta f_0(1+E)$ is zero free and $U$ is simply connected,
choose an ordinary branch $\ell_0$ of $\log f_0$. Then
\begin{equation}\label{eq:logprimitive}
 L=\ell_0+\sum_{n\ge1}\frac{(-1)^{n+1}}n E^n
 \quad\text{satisfies}\quad L'=F'/F.
\end{equation}
The positive support of $E$ makes the logarithmic family admissible.
The constant monomial $t^\beta$ contributes nothing to the derivative.
Equation~\eqref{eq:logprimitive} is a logarithmic primitive; identifying
it with a global inverse of a specified exponential requires additional
normalization, discussed in Section~\ref{sec:exponential}.

\subsection{Morera with a specified coefficient structure}
\begin{theorem}[Coefficientwise Morera theorem]\label{thm:morera}
Suppose a family $F=\sum_{\gamma\in S}f_\gamma t^\gamma$ has a
common well-ordered support and each $f_\gamma:U\to\C$ is
continuous. Define its contour integrals coefficientwise at ordinary
points. If its integral around every triangle with closed interior in
$U$ is zero, then every $f_\gamma$ is holomorphic, so $F$ has the
canonical evaluated extension of an element of $\HH_\Gamma(U)$.
\end{theorem}
\begin{proof}
Zero Hahn integrals have all coefficients zero. Each continuous
$f_\gamma$ therefore has zero integrals around triangles. Classical
Morera, or its proof by constructing a local path-independent
primitive on a disk, makes $f_\gamma$ holomorphic. Theorem
\ref{thm:Heval} supplies the extension.
\end{proof}
This does not claim that arbitrary fine-continuous surcomplex
functions satisfy Morera. Their coefficient functions and their
contours need not even be defined in this sense.

\subsection{Linear systems with infinitesimal deformations}
Let $\HH_{\Gamma,\ge0}(U)$ denote the subring of Hahn families with
nonnegative support, and let reduction mean the coefficient at zero.
Matrix and vector families use a common support obtained by a finite
union.

\begin{theorem}[Holomorphic linear systems over Hahn coefficients]\label{thm:ODE}
Let $U$ be simply connected, $a\in U$, and suppose
\[
 A\in M_d(\HH_{\Gamma,\ge0}(U)),\quad
 B\in\HH_{\Gamma,\ge0}(U)^d,\quad
 c\in (K_\Gamma\cap\Fin)^d.
\]
The initial-value problem
\[
 Y'=AY+B,\qquad Y(a)=c
\]
has a unique solution in $\HH_{\Gamma,\ge0}(U)^d$.
\end{theorem}
\begin{proof}
Let $M$ be the monoid generated by the union of all positive supports
in $A,B,c$. Write the unknown as $Y=\sum_{\lambda\in M}Y_\lambda
t^\lambda$. At exponent zero solve the ordinary holomorphic system
\[
 Y_0'=A_0Y_0+B_0,\qquad Y_0(a)=c_0.
\]
A holomorphic fundamental matrix exists on simply connected $U$;
variation of constants gives its unique global solution.
At a later exponent $\lambda$, solve
\[
 Y_\lambda'=A_0Y_\lambda+B_\lambda+
      \sum_{\substack{\alpha+\beta=\lambda\\\alpha>0}}
                    A_\alpha Y_\beta,
 \qquad Y_\lambda(a)=c_\lambda.
\]
The forcing is a finite sum of known holomorphic functions and
$\beta<\lambda$ whenever it occurs. Well-ordered recursion gives a
Hahn family with support in $M$. It satisfies the system coefficient
by coefficient and hence after evaluation. Comparing the first
nonzero coefficient of a difference proves uniqueness, by uniqueness
of the ordinary homogeneous initial-value problem.
\end{proof}

\begin{example}[Why the nonnegative-support hypothesis matters]
There is no nonzero $Y\in\HH_\Gamma(U)$ satisfying
$Y'=t^{-1}Y$ on any connected ordinary domain. Indeed,
$v_H(Y')\ge v_H(Y)$ whenever $Y'\ne0$, while
$v_H(t^{-1}Y)=v_H(Y)-1$, an impossibility. Nevertheless
$\sum_{n\ge0}t^{-n}(z-a)^n/n!$ is a perfectly valid local analytic
germ on a sufficiently small scale. A local solution need not extend
as a single coherent family over an ordinary-sized connected base.
\end{example}

\section{Formal residues, Lagrange inversion, and infinity}\label{sec:formalres}
The local algebra has its own residue calculus, independent of the
choice of a global coefficientwise contour. Let
$\SC((X))$ denote Laurent series with finitely many negative powers
and arbitrary surcomplex coefficients. They define meromorphic germs
on sufficiently small punctured neighborhoods.

\begin{theorem}[Formal residue calculus]\label{thm:formalres}
For $A(X)=\sum_{n\ge-N}a_nX^n$, define
$\Res_0 A(X)\,dX=a_{-1}$. Then:
\begin{enumerate}[label=(\roman*)]
\item a Laurent differential has a Laurent primitive if and only if
its residue is zero;
\item if $\psi(T)=T^mU(T)$, $m\ge1$, $U(0)\ne0$, then
\[
 \Res_{T=0}A(\psi(T))\psi'(T)\,dT
      =m\Res_{X=0}A(X)\,dX;
\]
\item for an invertible local coordinate, the residue is invariant.
\end{enumerate}
\end{theorem}
\begin{proof}
Formal differentiation never produces an $X^{-1}$ term. Conversely,
when $a_{-1}=0$, termwise integration gives
\[
 \sum_{n\ne-1}\frac{a_n}{n+1}X^{n+1},
\]
which is again a Laurent series with finite principal part.
For change of variable it suffices to check monomials. If $n\ne-1$,
$\psi^n\psi'=(\psi^{n+1})'/(n+1)$ has zero residue. For $n=-1$,
\[
 \frac{\psi'}\psi=\frac mT+\frac{U'}U,
\]
and the second term is a power series. This proves (ii); (iii) is
$m=1$. All substitutions are formal coefficientwise operations with
finite contributors, and their germ realizations are supplied by
Section~\ref{sec:germs}.
\end{proof}
Thus the local quotient of meromorphic one-forms by exact one-forms
is one-dimensional over $\SC$, represented by $dX/X$. This is the
precise local logarithmic obstruction.

\begin{corollary}[Lagrange--B\"urmann inversion]\label{cor:lagrange}
Let $\Phi(X)\in\SC[[X]]$ with $\Phi(0)\ne0$, let
$P(X)=X/\Phi(X)$, and let $Q$ be its formal inverse. For
$H\in\SC[[X]]$ and $n\ge1$,
\begin{equation}\label{eq:lagrange}
 [Y^n]H(Q(Y))=\frac1n[X^{n-1}]H'(X)\Phi(X)^n.
\end{equation}
The identity remains valid after all sufficiently small admissible
surcomplex substitutions.
\end{corollary}
\begin{proof}
Express the coefficient as a residue and substitute $Y=P(X)$:
\[
 [Y^n]H(Q(Y))=
 \Res_0 H(X)\frac{P'(X)}{P(X)^{n+1}}\,dX.
\]
The residue of $(H P^{-n})'$ is zero. Hence the last expression is
$\frac1n\Res_0 H'(X)P(X)^{-n}\,dX$, which is exactly
\eqref{eq:lagrange}. Realization follows from the local Taylor
calculus.
\end{proof}

\subsection{The algebraic sphere}
For $R(z)\in\SC(z)$, use $w=1/z$ at infinity and define
\[
 \Res_\infty R(z)\,dz
    :=\Res_{w=0}\left(-w^{-2}R(1/w)\right)\,dw.
\]

\begin{proposition}[Global rational residue theorem]\label{prop:rational}
On the algebraic projective line $\mathbb P^1(\SC)$,
\[
 \sum_{b\in\SC}\Res_b R(z)\,dz+\Res_\infty R(z)\,dz=0.
\]
Only finitely many terms occur. A rational function with no finite
poles is a polynomial. A rational function bounded by one fixed
surreal modulus wherever it is defined is constant.
\end{proposition}
\begin{proof}
Finite partial fractions over the algebraically closed field $\SC$
write $R$ as a polynomial plus terms $C_{bk}(z-b)^{-k}$. The
coefficient of $z^{-1}$ at infinity is $\sum_b C_{b1}$, while the
coordinate change contributes its negative as the residue at infinity.
These $C_{b1}$ are exactly the finite residues.

A nonconstant denominator in a reduced fraction has a root, so a
rational function without finite poles is a polynomial. A rational
function with a pole is unbounded on arbitrarily small increments
about that pole, by its leading Laurent term. A nonconstant polynomial
is unbounded at infinity: choose a positive real-surreal $z$ large
enough that its highest-degree term dominates the finite sum of all
lower-degree terms and exceeds the proposed bound. Such a $z$ exists
in an ordered real closed field. Therefore a globally bounded rational
function must be constant.
\end{proof}
This is an algebraic rational subclass, not a classification of every
fine-locally meromorphic function on a class-sized topological sphere.
The latter class permits independent local choices.

\section{Finite angles and nonunique global exponentials}\label{sec:exponential}
The canonical lift of the ordinary exponential is defined on all
finite surcomplex arguments. Write it as $\exp^\sharp$. Formal
addition and the ordinary exponential law give
\[
 \exp^\sharp(z+w)=\exp^\sharp(z)\exp^\sharp(w),
 \qquad z,w\in\Fin.
\]
Its derivative is itself and it has no zeros. On real finite arguments
it agrees with the canonical surreal exponential, denoted
$\exp_{\No}$; the latter is defined on all of $\No$ and obeys the
addition law and infinitesimal Taylor expansion \cite{vdDE}.

\subsection{The whole surcomplex unit circle already has finite angles}
\begin{theorem}[Finite-angle parametrization]\label{thm:angle}
Every $u\in\SC$ with $|u|=1$ has the form
\[
 u=\exp^\sharp(i\theta)
\]
for a finite real-surreal $\theta$. This angle is unique modulo
$2\pi\Z$. Thus, as an abelian group, the surcomplex unit circle is
\[
 (\Fin\cap\No)/(2\pi\Z).
\]
\end{theorem}
\begin{proof}
The real and imaginary parts of $u$ have modulus at most one, so $u$
is finite. Its standard part $u_0$ lies on the ordinary unit circle.
Choose an ordinary real angle $\theta_0$ with $e^{i\theta_0}=u_0$.
Then $h=u/u_0$ has the form $1+\delta$ with $\delta\in\mm$,
and $h\overline h=1$.

The formal infinitesimal logarithm satisfies
$\log_{\mm}(h)+\log_{\mm}(\overline h)=0$, while conjugation
commutes with its real Taylor coefficients. Therefore
$\log_{\mm}(h)$ is purely imaginary, say $i\epsilon$ with real
infinitesimal $\epsilon$. Exponentiation gives
$u=\exp^\sharp(i(\theta_0+\epsilon))$.

If two finite angles give the same value, subtract them and take
standard parts. Their standard difference is $2\pi k$ for an
ordinary integer $k$. The remaining difference is infinitesimal,
and the infinitesimal exponential is injective by its inverse
logarithm. Hence that difference is zero.
\end{proof}

\begin{corollary}[An obstruction to the classical exponential kernel]\label{cor:expobstruction}
There is no group homomorphism $E:(\SC,+)\to(\SC^*,\cdot)$
which simultaneously agrees with $\exp^\sharp$ on $\Fin$,
satisfies $|E(x+iy)|=\exp_{\No}(x)$, and has kernel exactly
$2\pi i\Z$.
\end{corollary}
\begin{proof}
The value $E(i\omega)$ has modulus one, so by
Theorem~\ref{thm:angle} it equals $\exp^\sharp(i\theta)=E(i\theta)$
for a finite $\theta$. Therefore $i(\omega-\theta)$ is in the
kernel. It is not an ordinary integral multiple of $2\pi i$.
\end{proof}
The obstruction is structural: it is not a failure of a particular
power-series prescription.

\subsection{Explicit entire analytic exponentials}
For a real surreal $y=\sum r_\gamma t^\gamma$, define the finite
truncation and one coefficient functional by
\[
 \pi_{\mathrm{fin}}(y)=\sum_{\gamma\ge0}r_\gamma t^\gamma,
 \qquad c_\omega(y)=r_{-1}.
\]
Both are additive. The first discards all purely infinite monomials;
the second extracts the coefficient of $\omega=t^{-1}$. On finite
$y$, $\pi_{\mathrm{fin}}(y)=y$ and $c_\omega(y)=0$.

\begin{theorem}[A family of global analytic exponentials]\label{thm:exponentials}
For every $\alpha\in\R$, the formula
\begin{equation}\label{eq:Ealpha}
 E_\alpha(x+iy)=\exp_{\No}(x)\,
 \exp^\sharp\!\left(i\bigl(\pi_{\mathrm{fin}}(y)
                            +\alpha c_\omega(y)\bigr)\right)
\end{equation}
defines a nonvanishing function on all of $\SC$ satisfying
\[
 E_\alpha(z+w)=E_\alpha(z)E_\alpha(w),\quad
 E_\alpha'=E_\alpha,\quad
 |E_\alpha(x+iy)|=\exp_{\No}(x).
\]
It agrees with $\exp_{\No}$ on the real-surreal axis and with the
canonical complex exponential on the entire finite plane. It is
locally Hahn analytic everywhere. Nevertheless
\[
 E_\alpha(i\omega)=e^{i\alpha},
\]
so these extensions are not uniquely determined by those properties.
\end{theorem}
\begin{proof}
The argument of $\exp^\sharp$ in \eqref{eq:Ealpha} is finite:
$c_\omega(y)$ and $\alpha$ are ordinary real numbers and the
finite truncation is finite. Additivity of the two maps and the two
exponential addition laws prove the multiplicative identity. The
phase has modulus one, proving the modulus formula and nonvanishing.
The restrictions to real and finite complex arguments follow from
the definitions.

For any infinitesimal $h=h_1+ih_2$, the two coefficient maps act as
$\pi_{\mathrm{fin}}(h_2)=h_2$ and $c_\omega(h_2)=0$. Hence
\[
 E_\alpha(z+h)=E_\alpha(z)\exp^\sharp(h)
       =E_\alpha(z)\Hs_{n\ge0}\frac{h^n}{n!}.
\]
This is an actual Taylor representation throughout $z+\mm$ and
proves local analyticity and the derivative formula. At $y=\omega$,
the finite truncation is zero and $c_\omega(y)=1$, giving the final
identity.
\end{proof}
For example, $E_0$ and $E_1$ already differ at $i\omega$ while
agreeing at every finite complex point and every real-surreal point.
For $E_0$, the kernel contains $iy$ for every purely infinite
real surreal $y$. These functions demonstrate why an
exponential addition law, the usual differential equation, conjugation
compatibility, and agreement on familiar arguments do not select a
unique global surcomplex exponential.

\begin{remark}[Two different notions of differentiation]
Throughout the article, every element of $\SC$ is a constant with
respect to the independent variable $z$. In particular,
$(t^\gamma)'=0$. This derivative must not be confused with a
nontrivial derivation of the surreal \emph{coefficient field}, such as
the derivations studied in \cite{BMderiv}. A field derivation can move
$t$; the analytic derivative used here does not.
\end{remark}

\section{Scale charts and infinite centers}\label{sec:scales}
The finite tube is a normalization, not a claim that infinite points
are inaccessible. Let $c\in\SC$, $r\in\SC^*$, and set
$z=c+rw$. A family $F\in\HH_\Gamma(U)$ defines
\[
 \mathcal F(z)=F\!\left(\frac{z-c}{r}\right)
       \quad\text{on }c+rU^\sharp.
\]
The derivative and differential transform as
\[
 \mathcal F^{(n)}(z)=r^{-n}F^{(n)}(w),\qquad dz=r\,dw.
\]
Define the integral along $c+r\gamma$ by pullback to $\gamma$.
Then preparation, zero counting, Cauchy's formula, and the residue
theorem carry over verbatim in normalized coordinates. In particular,
for a one-form the residue is invariant under the invertible change
of coordinate, as in Theorem~\ref{thm:formalres}.

\begin{example}[A contour at an infinite center and a very small scale]
Take $c=\omega^2$, $r=t^\omega$, and
\[
 \mathcal F(z)=\sin^\sharp\!\left(\frac{z-\omega^2}{t^\omega}\right)
 \quad\text{on }\omega^2+t^\omega\C^\sharp.
\]
Its zeros are exactly
\[
 z_k=\omega^2+k\pi t^\omega,\qquad k\in\Z,
\]
and each is simple. Let $\gamma$ be the ordinary circle
$|w|=3\pi/2$ and lift it by $z=\omega^2+t^\omega w$. There are
three enclosed normalized zeros, and
\[
 \frac1{2\pi i}\Hoint_{\omega^2+t^\omega\gamma}
        \frac{\mathcal F'(z)}{\mathcal F(z)}\,dz=3.
\]
The center is infinite and the radius is smaller than every finite
power of $t$, yet the normalized argument principle is exact.
\end{example}

Affine changes whose normalized form is an ordinary holomorphic map
plus a positive-support Hahn perturbation are covered by
Proposition~\ref{prop:Hcompose}. Completely arbitrary changes of scale
need not preserve one fixed globally supported family on a fixed
ordinary base. A global geometry must specify which charts and which
continuation data are allowed; it cannot be obtained merely by
forgetting this dependence.

\section{Worked examples and reproducible checks}\label{sec:examples}
\subsection{Preparation and splitting for \texorpdfstring{$z^2-te^z$}{z squared minus t exp z}}
Consider the Hahn family
\begin{equation}\label{eq:example1}
 F(z)=z^2-t\exp^\sharp(z)
 \quad\text{on }\C^\sharp.
\end{equation}
Its leading function is $z^2$. The zero-specialization theorem says
that it has exactly two zeros, counted with multiplicity, in the
entire finite plane, both infinitesimal.

Put $s=t^{1/2}$. The equation becomes
$z\exp_{\mm}(-z/2)=\pm s$. Lagrange inversion gives the inverse
series
\begin{equation}\label{eq:rho}
 \rho(s)=\sum_{n\ge1}
   \frac{n^{n-1}}{2^{n-1}n!}s^n
 =s+\frac{s^2}{2}+\frac{3s^3}{8}+\frac{s^4}{3}
   +\frac{125s^5}{384}+\frac{27s^6}{80}+\cdots.
\end{equation}
Indeed, take $\Phi(z)=e^{z/2}$ in
Corollary~\ref{cor:lagrange}. The two roots are exactly
\[
 z_+=\rho(t^{1/2}),\qquad z_-=\rho(-t^{1/2}).
\]
They are distinct, because their leading terms are $s$ and $-s$.
Every series in these formulas is a Hahn sum.

The prepared polynomial is
$W(z)=(z-z_+)(z-z_-)$, and its first coefficients are
\begin{align}\label{eq:Wexample}
 W(z)={}&z^2-t(1+z)
        -t^2\left(\frac12+\frac23z\right)\notag\\
       &-t^3\left(\frac{11}{24}+\frac{27}{40}z\right)
        +t^4\bigl(b_0(t)+b_1(t)z\bigr),
\end{align}
with $b_0,b_1\in\Q[[t]]$. These coefficients can also be obtained
without knowing the roots. In the preparation recursion, $u_0=1$,
$W_1=-(1+z)$, and
\[
 U_1=\frac{-e^z+1+z}{z^2}.
\]
At the next step $W_2=\operatorname{rem}_{<2}(-U_1W_1)$, yielding
$-\frac12-\frac23z$. The following step gives the displayed
$t^3$ coefficient.

For a small ordinary circle about zero, the logarithmic derivative
has expansion
\begin{equation}\label{eq:logderexample}
 \frac{F'}F=\frac2z-
  \sum_{n\ge1}\frac{t^n}{n}\frac d{dz}
           \left(\frac{e^{nz}}{z^{2n}}\right).
\end{equation}
Every positive-$t$ coefficient is a derivative of a meromorphic
function and has residue zero. Therefore the coefficientwise contour
integral is exactly $4\pi i$, agreeing with the two actual
surcomplex roots. This is a direct illustration of the full residue
and argument-principle theorems.

\subsection{A perturbation beyond every finite power}
Let
\[
 q=t^\omega,\qquad G(z)=z^2-t-q\exp^\sharp(z).
\]
Again the leading function is $z^2$, so there are exactly two finite
surcomplex zeros. Relative to the field with rational exponents,
$q$ is smaller than every positive coefficient scale. Start from
$a=\pm t^{1/2}$ and write $z=a+\delta$. The equation is
\[
 2a\delta+\delta^2=q\exp_{\mm}(a)\exp_{\mm}(\delta).
\]
Equivalently, $q=(2a\delta+\delta^2)e^{-a-\delta}$, whose
linear coefficient in $\delta$ is nonzero. Formal inversion over
$K_\Q$ therefore yields
\begin{equation}\label{eq:twoscale}
 z=a+q\frac{e^a}{2a}
       +q^2\frac{e^{2a}(2a-1)}{8a^3}+q^3R_a(q),
 \qquad R_a(q)\in K_\Q[[q]].
\end{equation}
Here $e^a$ denotes the infinitesimal Taylor exponential, and
Lemma~\ref{lem:ultrasmall} makes the entire $q$-series evaluable at
$q=t^\omega$. The first two correction valuations are
\[
 \omega-\frac12,\qquad 2\omega-\frac32.
\]
In particular, the first correction is smaller than every positive
finite power of $t$. No finite-order asymptotic expansion in $t$
alone would detect it, but it is an exact part of the surcomplex zero.
The prepared polynomial still has degree two, and the argument
integral around an ordinary circle is still exactly two after division
by $2\pi i$.

\subsection{What the accompanying script checks}
The archive contains a Python/SymPy script that uses exact rational
arithmetic to check four finite consequences of the formulas above:
the inverse-series equation for $\rho$ through degree ten; the
prepared polynomial coefficients through degree three in $t$; the
vanishing logarithmic-derivative residues for several positive
orders; and the two displayed $q$-correction coefficients.

The script operates in ordinary truncated formal series. It does not
attempt to store arbitrary surreal supports, decide well-ordering,
or verify the general theorems. Its role is to catch algebraic and
transcription errors in the worked examples. The proofs, rather than
numerical or symbolic experiments, establish all-order validity.

\section{Scope, limitations, and further development}\label{sec:scope}
The theory proved here has three connected layers. The class field
supplies arbitrary constants and scales. Local Hahn germs supply
analytic differential algebra. Coherent coefficient functions over an
ordinary connected base supply continuation, finite zero fibers,
and a contour calculus that measures their actual poles.

Several restrictions are mathematically essential. All supports and
all coefficient data in one construction are sets. A Hahn sum is not
a limit of its partial sums. Fine-topological local analyticity does
not imply coherence between monads. A finite tube omits monads over
the boundary of its classical base and omits infinite normalized
points. A single order bound on a function cannot control all its
Hahn coefficients. A contour integral is defined through classical
coefficient integration and admissible pullbacks, not through an
unproved surcomplex Riemann integration theory.

These restrictions still leave substantial scope. The preparation
recursion handles arbitrary ordered-rank infinitesimal perturbations,
not just one formal parameter or a discretely valued field. The
residue theorem resolves the difference between classical poles of
coefficient functions and displaced actual poles. The exact zero
specialization law explains why integer counts survive while roots
split through many scales. Scale charts move this machinery to
infinite centers and infinitesimal or infinite radii.

A more intrinsic global theory would need an admissible topology or
site of scale charts, with a rule excluding independent monadwise
choices while retaining useful changes of scale. It should distinguish
formal local completion from analytic continuation across a
classical-sized base. The linear equation $Y'=t^{-1}Y$ shows why
one fixed family ring cannot contain every analytically natural
solution at every scale. A several-variable version would require
uniformly supported holomorphic coefficient functions and a suitable
Weierstrass division in one distinguished variable. A functional
analytic version would need coefficient majorants compatible with
operator inversion and with contour deformation.

The results here do not settle these global design questions. They do
show that a substantial package of complex analysis survives over
$\No[i]$ once the analytic class, summation rule, and contour notion
are specified. The field's enormous supply of scales is not merely an
obstruction: it makes every formal local expansion realizable and
allows perturbations invisible at every finite asymptotic order to
be treated exactly.

\appendix
\section{Dependency and proof audit}\label{app:audit}
The logical order of the main arguments is as follows.
\begin{center}
\small
\begin{tabular}{@{}p{0.29\textwidth}p{0.62\textwidth}@{}}
\toprule
Result & Main inputs\\
\midrule
Local germ calculus & Normal forms, set-sized coefficient containment, support lemma, fine remainder estimates.\\
Coherent evaluation & $S+M(E)$ support control and classical Taylor identities.\\
Identity and maximum modulus & Monad-wide Taylor translation, classical coefficient identity, local dominant terms.\\
Weierstrass preparation & Positive support monoid, well-ordered recursion, holomorphic Taylor remainder division.\\
Zero specialization & Preparation, algebraic closedness of $K_\Gamma$, valuation comparison for polynomial roots.\\
Residue theorem & Preparation and division, finite partial fractions, classical contour integrals of coefficients.\\
Cauchy and argument principle & Residue theorem and local Taylor or Laurent expansions at actual poles.\\
Schwarz--Pick & Canonical lifting, divided differences, ordinary Schwarz--Pick and its equality case.\\
Global exponential examples & Finite Taylor exponential, canonical real-surreal exponential, additive normal-form truncations.\\
\bottomrule
\end{tabular}
\end{center}
No step uses convergence of an infinite sequence in the full fine
topology. No root is inferred solely from an infinitesimal residual:
roots are obtained from preparation in an algebraically closed
set-sized Hahn field. No division by a possibly vanishing leading
function is performed on its zero monad without first preparing the
denominator. No unrestricted interchange of infinite sums is used:
all such interchanges refer to the support lemma or to its
coefficient-field-small extension.

The foundational results imported rather than reproved are the
standard surreal normal-form and real-closedness theorems, the
Hahn--Neumann support lemma, algebraic closedness of Hahn fields with
divisible value group and algebraically closed coefficients, the
canonical surreal exponential's addition and infinitesimal Taylor
laws, and standard classical holomorphic-function theorems. The
surcomplex versions stated in this article are proved from these
inputs.

\clearpage
\phantomsection
\addcontentsline{toc}{section}{References}
\begin{thebibliography}{99}
\bibitem{Conway}
J. H. Conway, \emph{On Numbers and Games}, 2nd ed., A K Peters, 2001
(first edition, Academic Press, 1976).

\bibitem{Gonshor}
H. Gonshor, \emph{An Introduction to the Theory of Surreal Numbers},
London Mathematical Society Lecture Note Series 110, Cambridge
University Press, 1986.

\bibitem{vdDE}
L. van den Dries and P. Ehrlich, Fields of surreal numbers and
exponentiation, \emph{Fundamenta Mathematicae} \textbf{167} (2001),
173--188. In particular, Sections 1--2.
\url{https://www.impan.pl/shop/en/publication/transaction/download/product/89226}

\bibitem{Neumann}
B. H. Neumann, On ordered division rings,
\emph{Transactions of the American Mathematical Society}
\textbf{66} (1949), 202--252.
\url{https://doi.org/10.1090/S0002-9947-1949-0032593-5}

\bibitem{Poonen}
B. Poonen, Maximally complete fields,
\emph{L'Enseignement Math\'ematique} (2) \textbf{39} (1993), 87--106.
\url{https://math.mit.edu/~poonen/papers/amsval.pdf}

\bibitem{PoonenUnits}
B. Poonen, \emph{Units in Hahn--Mal'cev--Neumann rings}, manuscript,
14 January 2026.
\url{https://math.mit.edu/~poonen/papers/malcev.pdf}

\bibitem{BMgerms}
A. Berarducci and V. Mantova, Transseries as germs of surreal
functions, arXiv:1703.01995. In particular, Section 3 on summability,
analyticity, and composition.
\url{https://arxiv.org/abs/1703.01995}

\bibitem{BMderiv}
A. Berarducci and V. Mantova, Surreal numbers, derivations and
transseries, arXiv:1503.00315.
\url{https://arxiv.org/abs/1503.00315}

\bibitem{McMullen}
C. T. McMullen, \emph{Advanced Complex Analysis}, Harvard University
course notes, author-hosted version consulted 21 September 2026.
\url{https://people.math.harvard.edu/~ctm/papers/home/text/class/harvard/213a/course/course.pdf}

\bibitem{MS}
J. M\"uller and A. Strohmaier, The theory of Hahn-meromorphic
functions, a holomorphic Fredholm theorem, and its applications,
\emph{Analysis \& PDE} \textbf{7} (2014), 745--770.
\url{https://arxiv.org/abs/1205.0236}
\end{thebibliography}
\end{document}
